\chapter{Hedging a Structured-Products Book}\label{ch:s2:hedging-a-structured-products-book}

Retail investors buy autocallables in large size. The banks that sell them end up long volatility they must sell, short correlation and long
dividends, and their hedging moves those markets. The Bank for International Settlements noted in September 2026 that hedging the autocallables widely
sold to Korean retail investors can require large, destabilising trades near the barriers. On this chapter's synthetic book of \euro1 billion of
five-year autocallables, the bank is long \euro9.66 million of vega per volatility point and earns a margin of \euro8.31 million. Selling that vega into
a market that absorbs \euro10 million per point of concession lowers long-dated implied volatility by 0.97 points and costs the bank \euro4.67 million,
more than half its margin. The build is \texttt{firm.sphedge}.

\section{The autocallable book's exposures}

\begin{definition}[Exotic book exposure]\label{def:s2:hedging-a-structured-products-book:exposure}
An \emph{exotic book exposure}\index{exotic book exposure} is a sensitivity of a book of structured products to a market input that vanilla
instruments hedge only in part (volatility at long maturities, skew, dividends, correlation), measured by revaluing the whole book with the input
bumped.
\end{definition}

The bank has sold three notes (\cref{lst:s2:hedging-a-structured-products-book:book}): \euro400 million on index A alone, \euro300 million on index B
alone, and \euro300 million on the worse of the two. Each is Book~5's autocallable (chapter~18): five years, annual observations, autocall at 100\%, a
Phoenix coupon of 7\% (9\% for the worst-of) paid above a 70\% barrier with memory, and capital protected unless the index ends below 60\%. The
indices have 20\% volatility, a local-volatility skew that adds 10 points per 100\% fall, 3\% dividends and a correlation of 0.7. The single notes are
worth 99.12 per 100 and the worst-of 99.29: the bank's margin is \euro8.31 million.

\begin{center}
\small
\begin{tabular}{@{}lcccc@{}}
\toprule
issuer's P\&L for the bump (\euro{} millions) & single A & single B & worst-of & book\\
\midrule
volatility $+1$ point & 3.22 & 2.42 & 4.02 & 9.66\\
skew $+0.05$ & 1.47 & 1.10 & 1.86 & 4.42\\
dividend yield $+0.1$ point & 0.51 & 0.38 & 0.55 & 1.43\\
correlation $+0.05$ & & & $-1.05$ & $-1.05$\\
\bottomrule
\end{tabular}
\end{center}

The investor's note is, in large part, a short put that knocks in below 60\%. The bank holds the other side: it gains when volatility, downside skew or
dividends rise, all of which make the put dearer. The worst-of note adds a short position in correlation: when the indices move more alike, the worse
of the two is less likely to be bad, which helps the investor and costs the bank. These exposures are the mirror image of what the investor sold, and
the bank must sell them on.

\begin{dated}{2026-09}{Autocallables in Korea}\label{dat:s2:hedging-a-structured-products-book:korea}
The BIS Quarterly Review of September 2026 described autocallables as widely sold to Korean retail investors and noted that hedging them can require
large, destabilising trades near the barriers. Issuance referencing Samsung Electronics and SK Hynix, which together exceed half the KOSPI 200, rose
sharply in the first half of 2026.
\end{dated}

\section{Vega and skew}

The book's exposures move with the index (\cref{fig:s2:hedging-a-structured-products-book:profile}). At inception, the single note on A leaves the bank
long \euro3.22 million of vega, and its hedge holds \euro106 million of the index. If the index opened lower, the vega would rise to a peak of \euro4.45
million near 80\% of the strike, then fall and turn negative below about 52\%: once the protection is likely lost, the note is mostly a holding of
the index. The hedge grows as the index falls, to \euro334 million at 55\%, which means buying into a falling market. Below that it shrinks, to
\euro246 million at 40\%, which means selling as the market falls further. That second region is the destabilising trade near the barrier.
(\cref{lst:s2:hedging-a-structured-products-book:profile} computes both curves.)

\begin{omfigure}
\begin{tikzpicture}
\begin{axis}[width=11cm,height=5.4cm,axis y line*=right,axis x line=none,ylabel={\small index held as hedge (\euro{} m)},
  tick label style={font=\footnotesize},xmin=40,xmax=120,ymin=0,ymax=400,table/col sep=comma]
\addplot[omDef,thick,dashed] table[x=spot,y=hedge]{figdata/strategies-2/28-hedging-a-structured-products-book/profile.csv};\label{plot:s2:sphedge}
\end{axis}
\begin{axis}[width=11cm,height=5.4cm,axis y line*=left,xlabel={\small index level (\% of strike)},ylabel={\small issuer's vega (\euro{} m per point)},
  tick label style={font=\footnotesize},xmin=40,xmax=120,ymin=-3,ymax=5,grid=major,grid style={gray!25},table/col sep=comma,
  legend style={font=\scriptsize,at={(0.5,-0.3)},anchor=north},legend columns=2]
\addplot[omThm,thick,mark=*,mark size=1pt] table[x=spot,y=vega]{figdata/strategies-2/28-hedging-a-structured-products-book/profile.csv};
\addlegendentry{vega (left)}
\addlegendimage{/pgfplots/refstyle=plot:s2:sphedge}\addlegendentry{index held as delta hedge (right)}
\end{axis}
\end{tikzpicture}

\omcaption{The bank's vega and delta hedge for the \euro400 million single note on index A, at inception, if the index stood at each level relative to
the strike. Data: \texttt{s2\_sphedge.profile}.}\label{fig:s2:hedging-a-structured-products-book:profile}
\end{omfigure}

\section{Correlation and dividends}

\begin{definition}[Recycling trade]\label{def:s2:hedging-a-structured-products-book:recycling}
A \emph{recycling trade}\index{recycling trade} is a trade by which an issuer of structured products passes on to other investors the exposures its
book cannot hedge with vanilla instruments: selling long-dated volatility, selling dividends, buying correlation, usually at a concession.
\end{definition}

The bank's dividend exposure, \euro1.43 million per 0.1 point of dividend yield, is recycled by selling dividend futures or swaps, the supply chapter~26
modelled; Eurex describes its dividend derivatives as used to hedge dividends particularly for structured products. The short correlation, \euro1.05
million per 0.05 of correlation, is recycled by buying correlation: selling index volatility against single-name volatility, the other side of
chapter~2's dispersion trade. Each \omterm{def:s2:hedging-a-structured-products-book:recycling}{recycling trade} needs a buyer, and the buyers (hedge funds, dispersion traders, dividend investors) are paid for
taking the other side (\cref{fig:s2:hedging-a-structured-products-book:exposures}).

\begin{omfigure}
\begin{tikzpicture}
\begin{axis}[width=11cm,height=5.2cm,ybar,bar width=8pt,ylabel={\small issuer's P\&L for the bump (\euro{} m)},tick label style={font=\footnotesize},
  xtick=data,xticklabels={vega,skew,dividends,correlation},ymin=-1.5,ymax=4.5,grid=major,grid style={gray!25},table/col sep=comma,
  legend style={font=\scriptsize,at={(0.97,0.97)},anchor=north east}]
\addplot[fill=omThm!60,draw=omThm] table[x=k,y=n0]{figdata/strategies-2/28-hedging-a-structured-products-book/exposures.csv};
\addplot[fill=gray!40,draw=gray] table[x=k,y=n1]{figdata/strategies-2/28-hedging-a-structured-products-book/exposures.csv};
\addplot[fill=omDef!60,draw=omDef] table[x=k,y=n2]{figdata/strategies-2/28-hedging-a-structured-products-book/exposures.csv};
\legend{single A,single B,worst-of}
\end{axis}
\end{tikzpicture}

\omcaption{The synthetic book's exposures by note: the bank's P\&L for a 1 point rise in volatility, a 0.05 rise in skew, a 0.1 point rise in dividend
yield and a 0.05 rise in correlation. Data: \texttt{s2\_sphedge.table}.}\label{fig:s2:hedging-a-structured-products-book:exposures}
\end{omfigure}

\section{Recycling trades}

Selling \euro9.66 million of vega per point moves the market that absorbs it. If buyers take \euro$D$ of vega for each point of concession, long-dated
implied volatility falls by $9.66/D$ points; the bank pays half that move on average, and a buyer who holds until volatility returns earns all of it:

\begin{center}
\small
\begin{tabular}{@{}lccc@{}}
\toprule
depth (\euro{} m of vega per point) & 5 & 10 & 20\\
\midrule
move of long-dated implied volatility (points) & $-1.93$ & $-0.97$ & $-0.48$\\
bank's cost (\euro{} m) & 9.33 & 4.67 & 2.33\\
buyer's gain if volatility returns (\euro{} m) & 18.67 & 9.33 & 4.67\\
\bottomrule
\end{tabular}
\end{center}

The cost of recycling is of the order of the margin, \euro8.31 million. When many issuers sell the same exposures, depth is shared and the move is
larger. Supply that pushes long-dated volatility, dividends and correlation away from fair value is what the buyers of chapters~2, 3 and~26 are paid
for. The issuer's alternatives are to keep the risk, which its limits rarely allow, or to price the recycling cost into the notes.

\section{Strategy files}

\begin{strategyfile}{Vega recycling to hedge funds}\label{strat:s2:hedging-a-structured-products-book:vega}
\sfield{Who pays you, and why} The issuer, who must sell the long-dated vega its notes leave it with.
\sfield{Instruments and venues} Long-dated index options and variance swaps, over the counter.
\sfield{Signal} Long-dated implied volatility below its fair level after issuance waves.
\sfield{Sizing and execution} Buy what issuers sell; hold until volatility returns.
\sfield{Costs} Wide spreads on long-dated options; funding of the premium.
\sfield{How it dies} Issuance that never stops; a volatility regime that stays low.
\sfield{Horizon, capacity, infrastructure} Years.
\sfield{Backtest honestly} Long-dated implied volatility history as quoted, with its spreads.
\sfield{Sources} This chapter: 0.97 points of implied volatility for \euro9.66 million of vega at a depth of \euro10 million.
\end{strategyfile}

\begin{strategyfile}{Dividend recycling}\label{strat:s2:hedging-a-structured-products-book:dividends}
\sfield{Who pays you, and why} Issuers long dividends from their notes' hedges.
\sfield{Instruments and venues} Index dividend futures and swaps.
\sfield{Signal} Dividend futures below forecast dividends.
\sfield{Sizing and execution} Held to expiry.
\sfield{Costs} Futures spreads; recession cuts.
\sfield{How it dies} Recessions and dividend bans.
\sfield{Horizon, capacity, infrastructure} One to five years.
\sfield{Backtest honestly} Include recession years.
\sfield{Sources} Eurex factsheet; chapter~26; this chapter: \euro1.43 million per 0.1 point of dividend yield.
\end{strategyfile}

\begin{strategyfile}{Correlation recycling}\label{strat:s2:hedging-a-structured-products-book:corr}
\sfield{Who pays you, and why} Issuers of worst-of notes, short correlation.
\sfield{Instruments and venues} Correlation swaps; index against single-stock volatility (dispersion).
\sfield{Signal} Implied correlation against realised.
\sfield{Sizing and execution} Vega-weighted dispersion (chapter~2).
\sfield{Costs} Single-stock option spreads.
\sfield{How it dies} Crashes, when correlation jumps.
\sfield{Horizon, capacity, infrastructure} Months to years.
\sfield{Backtest honestly} Correlation in crashes.
\sfield{Sources} This chapter: $-$\euro1.05 million per 0.05 of correlation on a \euro300 million worst-of.
\end{strategyfile}

\begin{strategyfile}{Skew hedge programme}\label{strat:s2:hedging-a-structured-products-book:skew}
\sfield{Who pays you, and why} The issuer's own risk limits; the market's price for downside protection.
\sfield{Instruments and venues} Out-of-the-money index puts and put spreads.
\sfield{Signal} The book's skew exposure by maturity and level.
\sfield{Sizing and execution} Sell downside skew against the book's long position, more as the index nears the barriers.
\sfield{Costs} Put spreads; the destabilising hedge near barriers.
\sfield{How it dies} A crash through the barriers.
\sfield{Horizon, capacity, infrastructure} The notes' lives.
\sfield{Backtest honestly} Paths through the barriers, not only calm ones.
\sfield{Sources} BIS Quarterly Review, September 2026; this chapter: \euro4.42 million per 0.05 of skew.
\end{strategyfile}

\section{Tutorial: long what nobody wants}\label{tut:s2:hedging-a-structured-products-book:tut}

\begin{tutorial}
\textbf{Goal.} Price a book of autocallables, measure its exposures by revaluation, trace one note's vega and hedge against the index level, and cost
the recycling of its vega. \textbf{End state:} the three tables and two figures.
\begin{tutsteps}
\item \textbf{The book and its exposures}.
\omcode{code/firm/sphedge/firm_sphedge.py}{56}{79}{Notes on monthly paths; the issuer's P\&L for each bump.}\label{lst:s2:hedging-a-structured-products-book:book}
\item \textbf{Across the index level, and recycling}.
\omcode{code/firm/sphedge/firm_sphedge.py}{82}{102}{Vega and delta by level; the move and cost of selling vega.}\label{lst:s2:hedging-a-structured-products-book:profile}
\item \textbf{Run} \texttt{table()}, \texttt{margin()}, \texttt{profile()}, \texttt{recycling()} and \texttt{fig\_sphedge.py}.
\end{tutsteps}
\textbf{What to change next.} Age the book a year and recompute; add a third index and a worst-of on three; let depth fall when several issuers sell
at once.
\end{tutorial}

\section{Build: structured-products hedging}\label{bld:s2:hedging-a-structured-products-book:sphedge}

\begin{build}
\bfield{Purpose} Exposures of an autocallable book by revaluation, their profile across the index level, and the cost of recycling.
\bfield{Interface} \texttt{Note}, \texttt{BOOK}, \texttt{MarketState(\dots)}, \texttt{note\_value(note, mkt, paths, seed, spot)},
\texttt{book\_exposures}, \texttt{spot\_profile}, \texttt{recycling\_impact(vega, depth)}.
\bfield{Rules} Book 5's term sheet and cash flows; common random numbers for every bump; monthly steps (annual observations, knock-in at maturity).
\bfield{Acceptance tests} \texttt{code/firm/sphedge/tests/}: the book's signs; sensible note values; the worst-of gains from correlation; recycling
by hand.
\bfield{Stretch} Daily knock-in; ageing; several issuers sharing depth.
\end{build}

\begin{omsources}
\item Bank for International Settlements, \emph{BIS Quarterly Review}, September 2026, Box B, ``Short on memory? Leveraged ETFs and chip stock
volatility''.
\item Eurex, ``Dividend derivatives'', factsheet, accessed 25 September 2026.
\end{omsources}

\section{Exercises}

\begin{exercise}[$\star$]\label{exo:s2:hedging-a-structured-products-book:1}
A note is worth 99.12 per 100 and the bank sells \euro400 million of it at par. What is the margin?
\end{exercise}

\begin{exercise}[$\star$]\label{exo:s2:hedging-a-structured-products-book:2}
The book is long \euro9.66 million of vega per point and the market absorbs \euro5 million per point. How far does implied volatility fall, and
what does the bank pay?
\end{exercise}

\begin{exercise}[$\star$]\label{exo:s2:hedging-a-structured-products-book:3}
The delta hedge holds \euro106 million of the index at 100\% and \euro334 million at 55\%. What does the bank trade as the index falls between them?
\end{exercise}

\begin{exercise}[$\star\star$]\label{exo:s2:hedging-a-structured-products-book:4}
Why is the issuer of an autocallable long volatility?
\end{exercise}

\begin{exercise}[$\star\star$]\label{exo:s2:hedging-a-structured-products-book:5}
Why is the issuer of a worst-of note short correlation?
\end{exercise}

\begin{exercise}[$\star\star$]\label{exo:s2:hedging-a-structured-products-book:6}
Why can hedging near the barriers destabilise the market?
\end{exercise}

\begin{exercise}[$\star\star\star$]\label{exo:s2:hedging-a-structured-products-book:7}
\emph{Coding.} Read the vega profile. At what index level does the single note's vega change sign, and why?
\end{exercise}

\begin{exercise}[$\star\star\star$]\label{exo:s2:hedging-a-structured-products-book:8}
\emph{Find the flaw.} ``Our notes earn a 0.9\% margin with the risk hedged; recycling is just an operational detail.''
\end{exercise}

\section{Problem: Long What Nobody Wants}

\begin{problem}[{Weekend problem --- hedging a structured-products book}]\label{pb:s2:hedging-a-structured-products-book:1}
The chapter's synthetic book and the public record.

\medskip\noindent\textbf{Part I --- The book.}
\begin{enumerate}
\item Define an \omterm{def:s2:hedging-a-structured-products-book:exposure}{exotic book exposure}.
\item Describe the notes and the market.
\item What is the bank's margin?
\item What did the BIS observe in Korea?
\end{enumerate}

\medskip\noindent\textbf{Part II --- Exposures.}
\begin{enumerate}[resume]
\item Give the exposures table.
\item Why is the bank long volatility, skew and dividends?
\item Why is it short correlation?
\item How does vega change with the index level?
\end{enumerate}

\medskip\noindent\textbf{Part III --- Recycling.}
\begin{enumerate}[resume]
\item Define a \omterm{def:s2:hedging-a-structured-products-book:recycling}{recycling trade}.
\item How is each exposure recycled?
\item Give the recycling table.
\item Who is paid, and by whom?
\end{enumerate}

\medskip\noindent\textbf{Part IV --- The verdict.}
\begin{enumerate}[resume]
\item State the \emph{named result}: the book's net exposures and the implied-volatility move its recycling causes.
\item Why is the hedge destabilising below 55\%?
\item How does the recycling cost compare with the margin?
\item What happens when many issuers recycle at once?
\item Which earlier chapters are on the other side of this book?
\item Which strategy file is a trader's, not the issuer's?
\item How would you backtest a vega-recycling buyer?
\item In one sentence: what does the structured-products desk sell besides notes?
\end{enumerate}
\end{problem}

\section{Interview questions}

\begin{interviewq}[$\star$ \iqroles{trader}]\label{iq:s2:hedging-a-structured-products-book:1}
What exposures does a bank keep after selling autocallables?
\end{interviewq}

\begin{interviewq}[$\star\star$ \iqroles{researcher}]\label{iq:s2:hedging-a-structured-products-book:2}
How would you measure a structured-products book's correlation exposure?
\end{interviewq}

\begin{interviewq}[$\star\star$ \iqroles{trader}]\label{iq:s2:hedging-a-structured-products-book:3}
The index has fallen 40\% and your book is near its knock-in barriers. What do you do?
\end{interviewq}

\begin{interviewq}[$\star\star$ \iqroles{risk}]\label{iq:s2:hedging-a-structured-products-book:4}
What stress tests matter most for an autocallable book?
\end{interviewq}

\begin{interviewq}[$\star\star$ \iqroles{developer}]\label{iq:s2:hedging-a-structured-products-book:5}
Revaluing a book of ten thousand notes for twenty bumps takes too long. What would you do?
\end{interviewq}

\begin{interviewq}[$\star\star\star$ \iqroles{researcher}]\label{iq:s2:hedging-a-structured-products-book:6}
With a linear price concession, show that selling a quantity $V$ into depth $D$ costs $V^2/(2D)$, and that splitting the sale into $n$ equal parts
separated by full recovery costs $V^2/(2nD)$.
\end{interviewq}
